\documentclass[aspectratio=169,11pt]{beamer}
% ============================================================
% CONFIGURACIÓN
% ============================================================
\usetheme{Madrid}
\usecolortheme{default}
\usepackage[utf8]{inputenc}
\usepackage[T1]{fontenc}
\usepackage[spanish]{babel}
\usepackage{amsmath}
\usepackage{amssymb}
\usepackage{booktabs}
\usepackage{array}
\usepackage{graphicx}
\usepackage{xcolor}
\usepackage{listings}
\usepackage{tikz}
% ============================================================
% COLORES
% ============================================================
\definecolor{codebg}{RGB}{245,245,245}
\definecolor{codegreen}{RGB}{30,120,60}
\definecolor{codeblue}{RGB}{30,70,150}
\definecolor{codered}{RGB}{180,50,50}
% ============================================================
% PYTHON
% ============================================================
\lstdefinestyle{python}{
    language=Python,
    backgroundcolor=\color{codebg},
    basicstyle=\ttfamily\scriptsize,
    keywordstyle=\color{codeblue}\bfseries,
    commentstyle=\color{codegreen},
    stringstyle=\color{codered},
    numbers=left,
    numberstyle=\tiny\color{gray},
    numbersep=5pt,
    showstringspaces=false,
    breaklines=true,
    frame=single,
    tabsize=4
}
\lstset{style=python}
% ============================================================
% DATOS DE LA PRESENTACIÓN
% ============================================================
\title{EDA: Encontrar Patrones}
\subtitle{Data Detective}
\author{Análisis de Datos}
\institute{Universidad Marista}
\date{}
% ============================================================
% DOCUMENTO
% ============================================================
\begin{document}
% ============================================================
% PORTADA
% ============================================================
\begin{frame}
    \titlepage
\end{frame}
% ============================================================
% AGENDA
% ============================================================
\begin{frame}{Misión de hoy}
\begin{center}
\Large
\textbf{Convertir datos en evidencia}
\end{center}
\vspace{0.4cm}
Hoy aprenderemos a:
\begin{itemize}
    \item Conocer un dataset.
    \item Resumir información.
    \item Encontrar patrones.
    \item Detectar anomalías.
    \item Comparar grupos.
    \item Encontrar relaciones.
    \item Visualizar evidencia.
    \item Formular hipótesis.
    \item Defender conclusiones con datos.
\end{itemize}
\vspace{0.3cm}
\begin{alertblock}{Pregunta central}
¿Qué historia están contando nuestros datos?
\end{alertblock}
\end{frame}
% ============================================================
% TIMELINE
% ============================================================
\begin{frame}{Ruta de la investigación}
\begin{enumerate}
    \item Introducción a EDA
    \item Anatomía del dataset
    \item Estadística descriptiva
    \item Variables y distribuciones
    \item Outliers y anomalías
    \item Descanso
    \item Relaciones
    \item Visualización
    \item Data Detective
    \item Reto por equipos
    \item Conclusiones
\end{enumerate}
\end{frame}
% ============================================================
% BLOQUE 1
% ============================================================
\section{Introducción}
\begin{frame}{¿Qué es EDA?}
\textbf{EDA} significa:
\begin{center}
\Large
\textbf{Exploratory Data Analysis}
\end{center}
Es el proceso de:
\begin{itemize}
    \item explorar,
    \item resumir,
    \item visualizar,
    \item comparar,
    \item cuestionar
\end{itemize}
un conjunto de datos antes de construir modelos.
\begin{alertblock}{Idea importante}
EDA no busca confirmar rápidamente una idea.
Busca descubrir qué está pasando realmente.
\end{alertblock}
\end{frame}
% ------------------------------------------------------------
\begin{frame}{¿Por qué necesitamos EDA?}
Imaginemos un dataset:
\[
X = \{x_1,x_2,\ldots,x_n\}
\]
Antes de aplicar Machine Learning necesitamos saber:
\begin{itemize}
    \item ¿Qué representa cada variable?
    \item ¿Hay valores faltantes?
    \item ¿Hay errores?
    \item ¿Hay valores extremos?
    \item ¿Cómo se distribuyen los datos?
    \item ¿Existen grupos?
    \item ¿Existen relaciones?
\end{itemize}
\end{frame}
% ------------------------------------------------------------
\begin{frame}{EDA no es solamente hacer gráficas}
Un EDA completo combina:
\[
\boxed{
\text{Datos}
+
\text{Estadística}
+
\text{Visualización}
+
\text{Pensamiento crítico}
}
\]
\begin{itemize}
    \item Código para explorar.
    \item Estadística para resumir.
    \item Gráficas para observar.
    \item Contexto para interpretar.
\end{itemize}
\end{frame}
% ============================================================
% BLOQUE 2
% ============================================================
\section{Anatomía del dataset}
\begin{frame}{Anatomía de un dataset}
Un dataset puede representarse como una matriz:
\[
X =
\begin{bmatrix}
x_{11} & x_{12} & \cdots & x_{1p}\\
x_{21} & x_{22} & \cdots & x_{2p}\\
\vdots & \vdots & \ddots & \vdots\\
x_{n1} & x_{n2} & \cdots & x_{np}
\end{bmatrix}
\]
Donde:
\begin{itemize}
    \item $n$ = número de observaciones.
    \item $p$ = número de variables.
\end{itemize}
\end{frame}
% ------------------------------------------------------------
\begin{frame}{Filas vs columnas}
\begin{block}{Fila}
Representa normalmente una observación.
\end{block}
Ejemplo:
\begin{center}
\begin{tabular}{cccc}
\toprule
Edad & Carrera & Semestre & Promedio\\
\midrule
20 & MAC & 6 & 8.4\\
21 & MAC & 7 & 9.1\\
19 & Actuaría & 4 & 7.8\\
\bottomrule
\end{tabular}
\end{center}
Cada fila representa un estudiante.
\end{frame}
% ------------------------------------------------------------
\begin{frame}{Tipos de variables}
\textbf{Numéricas}
\begin{itemize}
    \item Edad
    \item Salario
    \item Promedio
    \item Temperatura
\end{itemize}
\vspace{0.3cm}
\textbf{Categóricas}
\begin{itemize}
    \item Carrera
    \item Ciudad
    \item Sexo
    \item Tipo de cliente
\end{itemize}
\vspace{0.3cm}
\textbf{Fecha / tiempo}
\begin{itemize}
    \item Fecha de compra
    \item Hora
    \item Año
\end{itemize}
\end{frame}
% ------------------------------------------------------------
\begin{frame}[fragile]{Primer contacto con Python}
\begin{lstlisting}
import pandas as pd
df = pd.read_csv("dataset.csv")
print(df.head())
\end{lstlisting}
Preguntas:
\begin{itemize}
    \item ¿Qué observamos?
    \item ¿Qué columnas existen?
    \item ¿Qué representa cada fila?
\end{itemize}
\end{frame}
% ------------------------------------------------------------
\begin{frame}[fragile]{Tamaño del dataset}
\begin{lstlisting}
df.shape
\end{lstlisting}
Ejemplo:
\begin{lstlisting}
(1000, 8)
\end{lstlisting}
Significa:
\[
1000 \text{ observaciones}
\]
y
\[
8 \text{ variables}
\]
\begin{alertblock}{Pregunta de detective}
¿1000 filas significa que tenemos 1000 personas?
\end{alertblock}
\end{frame}
% ------------------------------------------------------------
\begin{frame}[fragile]{Conocer las columnas}
\begin{lstlisting}
df.columns
\end{lstlisting}
También:
\begin{lstlisting}
df.info()
\end{lstlisting}
Nos permite conocer:
\begin{itemize}
    \item nombres,
    \item tipos de datos,
    \item valores no nulos,
    \item memoria utilizada.
\end{itemize}
\end{frame}
% ------------------------------------------------------------
\begin{frame}[fragile]{Valores faltantes}
\begin{lstlisting}
df.isnull().sum()
\end{lstlisting}
Ejemplo:
\begin{tabular}{cc}
\toprule
Variable & Faltantes\\
\midrule
Edad & 0\\
Salario & 12\\
Ciudad & 4\\
Promedio & 0\\
\bottomrule
\end{tabular}
\begin{alertblock}{Pregunta}
¿Por qué podrían faltar esos datos?
\end{alertblock}
\end{frame}
% ============================================================
% BLOQUE 3
% ============================================================
\section{Estadística descriptiva}
\begin{frame}{Estadística descriptiva}
La estadística descriptiva busca resumir los datos.
Tres preguntas:
\begin{enumerate}
    \item ¿Dónde están los datos?
    \item ¿Qué tan dispersos están?
    \item ¿Cómo se distribuyen?
\end{enumerate}
\end{frame}
% ------------------------------------------------------------
\begin{frame}{Media}
La media:
\[
\bar{x} =
\frac{1}{n}
\sum_{i=1}^{n}x_i
\]
Ejemplo:
\[
5,7,8,10
\]
\[
\bar{x}
=
\frac{5+7+8+10}{4}
=
7.5
\]
\begin{alertblock}{Cuidado}
La media puede verse afectada por valores extremos.
\end{alertblock}
\end{frame}
% ------------------------------------------------------------
\begin{frame}{Mediana}
La mediana es el valor central después de ordenar los datos.
Ejemplo:
\[
2,4,5,7,100
\]
La mediana es:
\[
5
\]
La media es:
\[
23.6
\]
\begin{block}{Insight}
La mediana puede representar mejor el centro cuando existen outliers.
\end{block}
\end{frame}
% ------------------------------------------------------------
\begin{frame}{Moda}
La moda es el valor más frecuente.
Ejemplo:
\[
1,2,2,2,3,4,5
\]
Entonces:
\[
\text{Moda}=2
\]
Es especialmente útil para variables categóricas.
\end{frame}
% ------------------------------------------------------------
\begin{frame}{Dispersión}
No basta con conocer el promedio.
Necesitamos saber qué tan separados están los datos.
Medidas:
\begin{itemize}
    \item Rango.
    \item Varianza.
    \item Desviación estándar.
    \item Rango intercuartílico.
\end{itemize}
\end{frame}
% ------------------------------------------------------------
\begin{frame}{Desviación estándar}
La desviación estándar mide qué tan alejados están los datos de la media.
\[
\sigma =
\sqrt{
\frac{1}{n}
\sum_{i=1}^{n}
(x_i-\bar{x})^2
}
\]
Una desviación estándar grande implica mayor dispersión.
\end{frame}
% ------------------------------------------------------------
\begin{frame}[fragile]{Pandas: resumen estadístico}
\begin{lstlisting}
df.describe()
\end{lstlisting}
Obtendremos:
\begin{itemize}
    \item count
    \item mean
    \item std
    \item min
    \item 25\%
    \item 50\%
    \item 75\%
    \item max
\end{itemize}
\end{frame}
% ============================================================
% BLOQUE 4
% ============================================================
\section{Variables y distribuciones}
\begin{frame}{¿Cómo se distribuyen los datos?}
Una distribución responde:
\begin{center}
\Large
¿Cómo se reparten los valores?
\end{center}
Podemos encontrar:
\begin{itemize}
    \item Distribuciones simétricas.
    \item Distribuciones sesgadas.
    \item Concentraciones.
    \item Múltiples grupos.
    \item Colas largas.
\end{itemize}
\end{frame}
% ------------------------------------------------------------
\begin{frame}{Histograma}
Un histograma divide los valores en intervalos llamados ``bins''.
Ejemplo conceptual:
\begin{center}
\begin{tikzpicture}[scale=0.7]
\draw[->] (0,0) -- (8,0);
\draw[->] (0,0) -- (0,4);
\draw[fill=gray!30] (0.5,0) rectangle (1.5,1);
\draw[fill=gray!30] (1.5,0) rectangle (2.5,2);
\draw[fill=gray!30] (2.5,0) rectangle (3.5,3.5);
\draw[fill=gray!30] (3.5,0) rectangle (4.5,3);
\draw[fill=gray!30] (4.5,0) rectangle (5.5,2);
\draw[fill=gray!30] (5.5,0) rectangle (6.5,1);
\end{tikzpicture}
\end{center}
\end{frame}
% ------------------------------------------------------------
\begin{frame}[fragile]{Crear un histograma}
\begin{lstlisting}
import matplotlib.pyplot as plt
df["edad"].hist(bins=10)
plt.xlabel("Edad")
plt.ylabel("Frecuencia")
plt.title("Distribucion de edades")
plt.show()
\end{lstlisting}
\end{frame}
% ------------------------------------------------------------
\begin{frame}{Preguntas frente a un histograma}
Nunca debemos solamente mirar la gráfica.
Preguntemos:
\begin{enumerate}
    \item ¿Dónde se concentran los datos?
    \item ¿Hay valores extremos?
    \item ¿Es simétrica?
    \item ¿Tiene una cola?
    \item ¿Hay más de un pico?
\end{enumerate}
\end{frame}
% ------------------------------------------------------------
\begin{frame}{Sesgo}
Una distribución puede estar sesgada.
\textbf{Sesgo positivo:}
\[
\text{Media} > \text{Mediana}
\]
\textbf{Sesgo negativo:}
\[
\text{Media} < \text{Mediana}
\]
Esto nos ayuda a interpretar correctamente la media.
\end{frame}
% ------------------------------------------------------------
\begin{frame}{Boxplot}
El boxplot resume:
\begin{itemize}
    \item mínimo,
    \item primer cuartil,
    \item mediana,
    \item tercer cuartil,
    \item máximo,
    \item posibles outliers.
\end{itemize}
El rango intercuartílico es:
\[
IQR = Q_3-Q_1
\]
\end{frame}
% ============================================================
% BLOQUE 5
% ============================================================
\section{Outliers}
\begin{frame}{¿Qué es un outlier?}
Un outlier es una observación que se encuentra muy alejada del comportamiento habitual.
Ejemplo:
\[
20,\;21,\;22,\;23,\;24,\;120
\]
El valor:
\[
120
\]
llama inmediatamente nuestra atención.
\end{frame}
% ------------------------------------------------------------
\begin{frame}{Método del IQR}
Definimos:
\[
IQR=Q_3-Q_1
\]
Límite inferior:
\[
Q_1-1.5(IQR)
\]
Límite superior:
\[
Q_3+1.5(IQR)
\]
Los valores fuera de estos límites pueden considerarse outliers.
\end{frame}
% ------------------------------------------------------------
\begin{frame}[fragile]{Detectar outliers con Python}
\begin{lstlisting}
Q1 = df["salario"].quantile(0.25)
Q3 = df["salario"].quantile(0.75)
IQR = Q3 - Q1
limite_inf = Q1 - 1.5 * IQR
limite_sup = Q3 + 1.5 * IQR
outliers = df[
    (df["salario"] < limite_inf) |
    (df["salario"] > limite_sup)
]
print(outliers)
\end{lstlisting}
\end{frame}
% ------------------------------------------------------------
\begin{frame}{Un outlier no significa error}
Un valor extremo puede ser:
\begin{itemize}
    \item Error de captura.
    \item Error de medición.
    \item Fraude.
    \item Evento excepcional.
    \item Cliente especial.
    \item Fenómeno real.
\end{itemize}
\begin{alertblock}{Regla}
No eliminamos un dato solamente porque sea extremo.
\end{alertblock}
\end{frame}
% ------------------------------------------------------------
\begin{frame}{DESCANSO}
\begin{center}
\Huge 10 minutos
\vspace{0.5cm}
\Large
Regresamos para investigar relaciones.
\end{center}
\end{frame}
% ============================================================
% BLOQUE 6
% ============================================================
\section{Relaciones}
\begin{frame}{Relaciones entre variables}
Ahora pasamos de:
\[
X
\]
a:
\[
X \longrightarrow Y
\]
Preguntamos:
\begin{center}
¿Existe una relación entre dos variables?
\end{center}
Ejemplos:
\begin{itemize}
    \item Horas de estudio $\rightarrow$ promedio.
    \item Publicidad $\rightarrow$ ventas.
    \item Edad $\rightarrow$ salario.
    \item Temperatura $\rightarrow$ consumo eléctrico.
\end{itemize}
\end{frame}
% ------------------------------------------------------------
\begin{frame}{Covarianza}
La covarianza mide cómo cambian dos variables conjuntamente.
\[
Cov(X,Y)
=
E[(X-\mu_X)(Y-\mu_Y)]
\]
Interpretación:
\begin{itemize}
    \item Positiva: tienden a crecer juntas.
    \item Negativa: una crece mientras otra disminuye.
    \item Cercana a cero: poca relación lineal.
\end{itemize}
\end{frame}
% ------------------------------------------------------------
\begin{frame}{Correlación}
El coeficiente de Pearson:
\[
r =
\frac{Cov(X,Y)}
{\sigma_X\sigma_Y}
\]
Su rango es:
\[
-1\leq r\leq1
\]
\begin{itemize}
    \item $r\approx1$: relación positiva fuerte.
    \item $r\approx-1$: relación negativa fuerte.
    \item $r\approx0$: poca relación lineal.
\end{itemize}
\end{frame}
% ------------------------------------------------------------
\begin{frame}{Correlación no implica causalidad}
Si encontramos:
\[
X \leftrightarrow Y
\]
no significa automáticamente:
\[
X \rightarrow Y
\]
Puede existir:
\begin{itemize}
    \item Una tercera variable.
    \item Una relación accidental.
    \item Sesgo de selección.
    \item Datos insuficientes.
\end{itemize}
\begin{alertblock}{Data Detective}
Una correlación es una pista, no una sentencia.
\end{alertblock}
\end{frame}
% ------------------------------------------------------------
\begin{frame}[fragile]{Matriz de correlación}
\begin{lstlisting}
corr = df.corr(numeric_only=True)
print(corr)
\end{lstlisting}
Podemos buscar:
\[
|r| \approx 1
\]
para encontrar relaciones lineales fuertes.
\end{frame}
% ============================================================
% BLOQUE 7
% ============================================================
\section{Visualización}
\begin{frame}{Visualizar es construir evidencia}
Una gráfica debe responder una pregunta.
No:
\begin{center}
``Hagamos una gráfica.''
\end{center}
Sino:
\begin{center}
``Quiero comprobar qué ocurre con X cuando Y cambia.''
\end{center}
\end{frame}
% ------------------------------------------------------------
\begin{frame}{¿Qué gráfica usar?}
\begin{tabular}{ll}
\toprule
Pregunta & Gráfica\\
\midrule
¿Cómo se distribuye? & Histograma\\
¿Hay outliers? & Boxplot\\
¿Dos variables se relacionan? & Scatterplot\\
¿Cómo cambia en el tiempo? & Línea\\
¿Comparar categorías? & Barras\\
¿Muchas correlaciones? & Heatmap\\
\bottomrule
\end{tabular}
\end{frame}
% ------------------------------------------------------------
\begin{frame}[fragile]{Scatterplot}
\begin{lstlisting}
plt.scatter(
    df["horas_estudio"],
    df["promedio"]
)
plt.xlabel("Horas de estudio")
plt.ylabel("Promedio")
plt.show()
\end{lstlisting}
Pregunta:
\begin{center}
¿Los estudiantes que estudian más obtienen mejores resultados?
\end{center}
\end{frame}
% ------------------------------------------------------------
\begin{frame}{Una buena visualización}
Debe tener:
\begin{itemize}
    \item Título.
    \item Ejes.
    \item Unidades.
    \item Escala adecuada.
    \item Leyenda cuando sea necesaria.
    \item Contexto.
\end{itemize}
\begin{alertblock}{Principio}
La gráfica debe hacer evidente el patrón, no esconderlo.
\end{alertblock}
\end{frame}
% ============================================================
% BLOQUE 8
% ============================================================
\section{Data Detective}
\begin{frame}{Caso: El misterio del rendimiento}
Tenemos información de estudiantes:
\begin{itemize}
    \item Edad.
    \item Horas de estudio.
    \item Asistencia.
    \item Promedio.
    \item Carrera.
    \item Horas de sueño.
\end{itemize}
\begin{block}{Misión}
Descubrir qué factores parecen estar relacionados con el rendimiento académico.
\end{block}
\end{frame}
% ------------------------------------------------------------
\begin{frame}[fragile]{Pista 1: Conocer la escena}
Primero:
\begin{lstlisting}
df.head()
df.shape
df.info()
df.describe()
\end{lstlisting}
Preguntas:
\begin{enumerate}
    \item ¿Cuántas observaciones existen?
    \item ¿Qué variables tenemos?
    \item ¿Hay datos faltantes?
    \item ¿Hay valores imposibles?
\end{enumerate}
\end{frame}
% ------------------------------------------------------------
\begin{frame}{Pista 2: Distribuciones}
Analizar:
\begin{itemize}
    \item Edad.
    \item Horas de estudio.
    \item Promedio.
    \item Horas de sueño.
\end{itemize}
Utilizar:
\begin{itemize}
    \item Histogramas.
    \item Boxplots.
    \item Estadística descriptiva.
\end{itemize}
\end{frame}
% ------------------------------------------------------------
\begin{frame}{Pista 3: Relaciones}
Investigar:
\[
HorasEstudio \leftrightarrow Promedio
\]
\[
Asistencia \leftrightarrow Promedio
\]
\[
Sueño \leftrightarrow Promedio
\]
Utilizar:
\begin{itemize}
    \item Scatterplots.
    \item Correlaciones.
\end{itemize}
\end{frame}
% ------------------------------------------------------------
\begin{frame}{Pista 4: El sospechoso}
Encontrar:
\begin{itemize}
    \item El estudiante con menor promedio.
    \item El estudiante con mayor promedio.
    \item El estudiante con más horas de estudio.
    \item Valores extremos.
\end{itemize}
Después preguntar:
\begin{center}
¿Son errores o casos reales?
\end{center}
\end{frame}
% ------------------------------------------------------------
\begin{frame}{Pista 5: La hipótesis}
Cada equipo debe formular una hipótesis.
Ejemplo:
\begin{quote}
``Los estudiantes que estudian más horas tienden a obtener mejores promedios.''
\end{quote}
Después:
\[
\text{Hipótesis}
\rightarrow
\text{Evidencia}
\rightarrow
\text{Conclusión}
\]
\end{frame}
% ============================================================
% BLOQUE 9
% ============================================================
\section{Reto}
\begin{frame}{Reto por equipos}
\begin{center}
\Huge
\textbf{DATA DETECTIVE}
\end{center}
\vspace{0.4cm}
Cada equipo recibirá un dataset.
Su misión:
\begin{enumerate}
    \item Explorar.
    \item Encontrar un patrón.
    \item Encontrar una anomalía.
    \item Encontrar una relación.
    \item Crear evidencia visual.
    \item Presentar una conclusión.
\end{enumerate}
\end{frame}
% ------------------------------------------------------------
\begin{frame}{Reglas del reto}
Cada equipo debe entregar:
\begin{enumerate}
    \item Una pregunta.
    \item Una hipótesis.
    \item Una estadística.
    \item Una gráfica.
    \item Un hallazgo.
    \item Una anomalía.
    \item Una conclusión.
\end{enumerate}
\begin{alertblock}{Importante}
No gana quien haga más gráficas.
Gana quien encuentre una historia interesante y pueda demostrarla.
\end{alertblock}
\end{frame}
% ------------------------------------------------------------
\begin{frame}{Plantilla de investigación}
\begin{block}{Pregunta}
¿Qué queremos descubrir?
\end{block}
\begin{block}{Hipótesis}
¿Qué creemos que ocurre?
\end{block}
\begin{block}{Evidencia}
¿Qué muestran los datos?
\end{block}
\begin{block}{Conclusión}
¿Qué podemos afirmar?
\end{block}
\begin{block}{Limitación}
¿Qué todavía NO podemos afirmar?
\end{block}
\end{frame}
% ============================================================
% BLOQUE 10
% ============================================================
\section{Conclusiones}
\begin{frame}{¿Qué aprendimos?}
EDA nos permite pasar de:
\begin{center}
\Large
\textbf{Datos}
\end{center}
a:
\[
\boxed{
\text{Datos}
\rightarrow
\text{Patrones}
\rightarrow
\text{Evidencia}
\rightarrow
\text{Hipótesis}
}
\]
\end{frame}
% ------------------------------------------------------------
\begin{frame}{El proceso completo}
\begin{enumerate}
    \item Cargar los datos.
    \item Conocer las variables.
    \item Revisar calidad.
    \item Describir estadísticamente.
    \item Analizar distribuciones.
    \item Detectar outliers.
    \item Analizar relaciones.
    \item Visualizar.
    \item Formular hipótesis.
    \item Comunicar hallazgos.
\end{enumerate}
\end{frame}
% ------------------------------------------------------------
\begin{frame}{Checklist del Data Scientist}
Antes de modelar:
\begin{itemize}
    \item[$\square$] ¿Conozco mi dataset?
    \item[$\square$] ¿Sé qué representa cada variable?
    \item[$\square$] ¿Revisé valores faltantes?
    \item[$\square$] ¿Revisé duplicados?
    \item[$\square$] ¿Encontré valores extraños?
    \item[$\square$] ¿Conozco las distribuciones?
    \item[$\square$] ¿Busqué relaciones?
    \item[$\square$] ¿Visualicé los patrones?
    \item[$\square$] ¿Tengo hipótesis?
    \item[$\square$] ¿Puedo defender mis conclusiones?
\end{itemize}
\end{frame}
% ------------------------------------------------------------
\begin{frame}{La idea que deben recordar}
\begin{center}
\Huge
\textbf{No analices datos\\
para hacer gráficas.}
\vspace{0.5cm}
\Large
Analiza datos para\\
\textbf{hacer preguntas mejores.}
\end{center}
\end{frame}
% ------------------------------------------------------------
\begin{frame}{Próxima misión}
\begin{center}
\Huge
\textbf{EDA $\rightarrow$ Machine Learning}
\vspace{0.5cm}
\Large
Una vez que conocemos nuestros datos...
\vspace{0.3cm}
\[
\text{¿Podemos hacer predicciones?}
\]
\end{center}
\end{frame}
% ------------------------------------------------------------
\begin{frame}
\begin{center}
\Huge
\textbf{Preguntas}
\vspace{1cm}
\Large
¿Quién encontró el patrón más interesante?
\end{center}
\end{frame}
\end{document}
